
\section*{Introdução}
\addcontentsline{toc}{section}{Introdução}
\label{sec:introducao}
% (==>) Este é o fenômeno; esta é a questão; há evidência científica suficiente para justificar investigá-la

% ECONOMIA --> MEIO AMBIENTE --> ECONOMIA
% Retroalimentação: decisões econômicas afetam o estoque ambiental, que pode alterar as condições produtivas e os retornos econômicos.

% HIPÓTESES
% ---
% HG: ECONOMIA E MEIO AMBIENTE FORMAM UM SISTEMA INTERDEPENDENTE
% Retornos + instituições --> conversão/conservação --> estoque de vegetação
% Estoque de vegetação × choques climáticos --> produtividade --> retornos econômicos

% H1: RETORNOS E INSTITUIÇÕES AFETAM CONVERSÃO/CONSERVAÇÃO
% Retorno econômico + instituições --> decisão de uso da terra --> estoque de vegetação
% Heterogeneidade territorial/biomas pode alterar incentivos e respostas dos agentes.

% H2: VEGETAÇÃO MODIFICA A RESPOSTA DA PRODUTIVIDADE A CHOQUES CLIMÁTICOS
% Choque climático (exógeno) × estoque de vegetação --> produtividade
% Heterogeneidade entre biomas pode alterar a magnitude dessa relação.

% EXTENSÃO POSSÍVEL:
% Informação conjunta de H1 e H2 pode subsidiar desenho e focalização territorial de políticas de conservação. Não constitui, por ora, hipótese independente.

% EIXOS CONCEITUAIS QUE CONECTAM O ARTIGO À TESE:
% ---
% ENDOGENEIDADE AMBIENTAL
% O estado ambiental não é apenas resultado da atividade econômica:
% ele também pode alterar resultados produtivos e incentivos econômicos.

% GOVERNANÇA TERRITORIAL
% Instituições, enforcement e incentivos operam sobre agentes e territórios heterogêneos, produzindo respostas distintas de conversão e conservação.

% INTERTEMPORALIDADE
% Decisões presentes de uso da terra alteram o estoque ambiental relevante para períodos posteriores e podem, por essa via, modificar condições futuras de produção. Por ora, é mecanismo conceitual, não hipótese testada diretamente.

% ------------------------------------
%%% (I) CONTEXTO
% ------------------------------------
% CITAÇÃO1–6 — H1: economia da conversão, retornos, instituições e enforcement.
% CITAÇÃO7–11 — H2: vegetação, clima, produtividade e identificação econométrica.
% CITAÇÃO12–14 — integração e Mato Grosso: feedback economia-ambiente, heterogeneidade e contexto empírico.

%% Conflito: preservar ou converter.
A expansão agropecuária costuma ser descrita como a escolha entre conservar a vegetação nativa para fins de sustentabilidade ou converter a terra a fim de obter maior retorno econômico {CITAÇÃO1}. Essa representação, contudo, simplifica uma decisão que não depende apenas da aptidão agronômica da terra ou da rentabilidade esperada da produção. A conversão da vegetação é uma decisão privada realizada sob preços, custos de produção, restrições de crédito, expectativas, liquidez, direitos de propriedade, regras ambientais, \textit{enforcement} e probabilidade  de sanções pelo uso irregular ou insustentável {CITAÇÃO2}.

Assim, territórios com expectativas de retornos agrícolas semelhantes podem apresentar distintos níveis de conversão quando diferem na configuração e na efetividade das instituições que regulam o uso da terra. Isto posto, uma mesma regra formal pode produzir respostas diferentes conforme a atividade econômica, o custo de oportunidade da terra e as condições locais de fiscalização {CITAÇÃO3}. Sob essa perspectiva, a conversão ou conservação da vegetação pode ser tratada como resultado endógeno da interação entre os retornos econômicos dos usos alternativos da terra e os incentivos, custos e riscos associados ao ambiente institucional {CITAÇÃO4}.

% Ponte para H1
Nessa toada, mudanças nos preços relativos, na rentabilidade agrícola ou nas condições de acesso aos mercados podem aumentar o retorno esperado da conversão, enquanto fiscalização, sanções, restrições legais ou outros mecanismos institucionais podem alterar a resposta dos agentes a esses incentivos e a intensidade dessa resposta também pode variar conforme o contexto econômico e as características dos biomas {CITAÇÃO5}. Portanto, compreender a evolução do estoque de vegetação requer analisar não apenas os retornos econômicos associados ao uso da terra, mas também as instituições sob as quais esses retornos são percebidos e transformados em decisões de conversão ou conservação {CITAÇÃO6}.

% Ponte para H2
Entretanto, o estoque ambiental afetado por essas decisões não é alheio ao processo produtivo. A vegetação nativa interfere diretamente nas condições biofísicas que sustentam a produção e, por essa via, tende a afetar os resultados econômicos, por estar associada a processos relacionados à umidade, precipitação e temperatura, que, por conseguinte, impactam a dinâmica físico-química e biológica do solo {CITAÇÃO7}. Isso ocorre, sobretudo, porque o estoque local de vegetação pode modificar as condições ecológicas do território e, consequentemente, a resposta da agricultura a choques climáticos {CITAÇÃO8}.

Essa distinção é particularmente importante porque variações climáticas de curto prazo podem fornecer uma fonte de variação plausivelmente exógena para investigar se a resposta da produtividade aos choques depende do estoque de vegetação previamente existente {CITAÇÃO9}. Se áreas com diferentes estoques prévios de vegetação respondem de maneira distinta a choques comparáveis de precipitação, temperatura ou déficit hídrico, então o estoque ambiental pode modificar o efeito desses choques sobre a produtividade agrícola {CITAÇÃO10}. A questão deixa, portanto, de ser apenas se a vegetação está correlacionada com maior ou menor produtividade e passa a ser se ela altera a sensibilidade da produção a perturbações externas, e se essa relação apresenta magnitude economicamente relevante e heterogeneidade entre biomas {CITAÇÃO11}.

% Justificativa para a integração das duas hipóteses
Os dois mecanismos formam, em conjunto, uma relação de codependência entre economia e meio ambiente. De um lado, retornos econômicos e instituições condicionam a decisão de converter ou conservar a vegetação e, consequentemente, a evolução do estoque ambiental. De outro, esse estoque pode modificar as condições biofísicas da produção e a resposta da atividade agrícola a choques climáticos, afetando a produtividade e os retornos econômicos associados ao uso da terra {CITAÇÃO12}. Assim, decisões econômicas alteram o estado ambiental, enquanto o estado ambiental pode alterar os resultados econômicos que fundamentam decisões subsequentes.

Enfatize-se que essa retroalimentação entre meio ambiente e produtividade pode ser especialmente relevante quando seus componentes são espacialmente heterogêneos. Os retornos da conversão, os custos institucionais enfrentados pelos agentes e a contribuição produtiva da vegetação podem variar entre atividades, territórios e biomas {CITAÇÃO13}. Consequentemente, políticas que tratam a conservação apenas como restrição à atividade econômica, ou que consideram separadamente os incentivos à conversão e os efeitos produtivos do estoque ambiental, podem deixar de capturar parte relevante da relação entre produção e conservação. Surge, então, uma questão econômica central: em que medida a interação entre incentivos privados, instituições e contribuição produtiva do estoque ambiental pode informar o desenho e a focalização territorial das políticas de conservação?

%% Mato Grosso neste contexto.
Nesse contexto, Mato Grosso constitui um caso empírico especialmente informativo para a investigação dessas relações. O estado combina produção agrícola e pecuária em larga escala com três biomas — Cerrado, Pantanal e Amazônia — que apresentam peculiaridades ambientais e econômicas submetidas a usos ora predatórios, ora racionais e sustentáveis. Além disso, observa-se, dentro de um mesmo território estadual, diferenças relevantes de vegetação, clima, disponibilidade hídrica, sistemas produtivos e custo de oportunidade da terra. Esses três biomas representam, portanto, não apenas formações ecológicas distintas, mas também contextos econômicos diferenciados, nos quais podem variar tanto a contribuição do estoque de vegetação para as condições produtivas quanto os incentivos privados associados à conversão ou à conservação da terra {CITAÇÃO14}.

Sendo assim, essa heterogeneidade permite investigar empiricamente os dois mecanismos discutidos anteriormente. Do lado da decisão de uso da terra, ela possibilita examinar se retornos econômicos e instituições produzem respostas distintas de conversão entre atividades e biomas. Nesse caso, não basta identificar se uma regra ambiental ou uma alteração na rentabilidade afeta, em média, o desmatamento ou a conversão; importa verificar se a intensidade dessas respostas depende das condições econômicas e territoriais sobre as quais os incentivos incidem. Assim, a heterogeneidade entre Amazônia, Cerrado e Pantanal interessa à pesquisa não apenas como controle espacial, mas como parte dos próprios mecanismos econômicos em análise.

Do lado da produção, importa saber se o estoque de vegetação modifica a resposta da produtividade aos choques climáticos e se essa relação varia entre biomas. O objetivo não é apenas verificar se a vegetação está associada ao nível médio de produtividade, mas também se ela altera a intensidade com que condições climáticas adversas afetam os resultados produtivos.

% ------------------------------------
%%% (II) PROBLEMA
% ------------------------------------

Diante desse quadro, o problema central deste artigo consiste em compreender como incentivos econômicos e institucionais moldam a conversão ou manutenção da vegetação e, simultaneamente, em que medida o estoque de vegetação previamente existente condiciona a resposta da produtividade agrícola a choques climáticos em ambientes ecologicamente e economicamente heterogêneos. A articulação desses dois mecanismos permite examinar uma relação de interdependência na qual decisões econômicas alteram o estado ambiental, enquanto esse estado ambiental, por sua vez, pode influenciar resultados produtivos que retroalimentam decisões subsequentes.

O \textbf{problema} do artigo pode ser sintetizado no seguinte questionamento: “Como retornos econômicos e instituições condicionam as decisões de conversão ou manutenção da vegetação, e em que medida o estoque de vegetação nativa altera a resposta da produtividade agrícola a choques climáticos em contextos ecologicamente e economicamente heterogêneos?”

Essa pergunta desdobra-se em duas dimensões empiricamente relacionadas: (i) como retornos econômicos, custos, riscos e incentivos institucionais se relacionam com as decisões de conversão ou manutenção da vegetação; e (ii) se o estoque de vegetação previamente existente modifica, em magnitude economicamente relevante, o efeito de choques climáticos sobre a produtividade agrícola, e se essa relação varia entre biomas. Em conjunto, essas duas dimensões buscam esclarecer como incentivos privados, instituições e condições ambientais interagem na relação entre produção e conservação, oferecendo subsídios para o desenho e a focalização territorial de políticas ambientais.

% ------------------------------------
%%% (III) OBJETIVOS E JUSTIFICATIVA
% ------------------------------------

O \textbf{objetivo geral} deste artigo é investigar como retornos econômicos e instituições se relacionam com as decisões de conversão ou manutenção da vegetação e se o estoque de vegetação previamente existente modifica a resposta da produtividade agrícola a choques climáticos em contextos territorialmente heterogêneos. De forma específica, o artigo busca: (i) examinar como retornos econômicos e incentivos institucionais se relacionam com as decisões de conversão ou manutenção da vegetação; (ii) verificar se estoques de vegetação alteram a resposta da produtividade agrícola a choques climáticos; e (iii) avaliar se esses mecanismos apresentam heterogeneidade entre biomas.

A relevância deste trabalho decorre, em termos teóricos, de tratar o meio ambiente não apenas como restrição ou externalidade, mas como componente endógeno do processo econômico, capaz de afetar resultados produtivos e, simultaneamente, ser transformado por decisões econômicas e institucionais. Essa relação envolve também uma dimensão territorial, na medida em que incentivos e instituições podem gerar respostas econômicas e ecológicas distintas conforme as condições locais, e uma dimensão intertemporal, pois decisões presentes de uso da terra alteram estoques ambientais que condicionam oportunidades e riscos econômicos futuros. Em termos empíricos e aplicados, compreender essas relações pode: (a) orientar decisões de firmas e produtores sobre uso da terra, risco e adaptação; (b) revelar repercussões para consumidores e mercados por meio da oferta, dos preços e da exposição ao risco; e (c) subsidiar governos e organizações do terceiro setor responsáveis por fiscalização, conservação e alocação de recursos ambientais.
% ------------------------------------
%%% (IV) HIPÓTESES
% ------------------------------------

Dito isso, a análise parte da seguinte \textbf{hipótese geral (HG):} retornos econômicos e instituições condicionam decisões de uso da terra que alteram o estoque ambiental, enquanto esse estoque pode, por sua vez, condicionar resultados produtivos relevantes para decisões econômicas subsequentes. Essa hipótese geral é examinada por meio de duas \textbf{hipóteses específicas}.

\medskip

\noindent{\textbf{(H1) Retornos, instituições e conversão:}}

\noindent{\emph{A resposta da conversão da vegetação aos retornos econômicos dos usos alternativos da terra é condicionada pelos incentivos, custos e riscos associados ao ambiente institucional, podendo variar entre atividades, territórios e biomas.}}

\medskip

\noindent{\textbf{(H2) Vegetação, clima e produtividade.}}

\noindent{\emph{O estoque de vegetação previamente existente modifica, em magnitude economicamente relevante, a resposta da produtividade agrícola aos choques climáticos, podendo essa relação variar entre biomas.}}

\medskip

% ------------------------------------
%%% (V) METODOLOGIA
% ------------------------------------

% Falta melhorar: detalhar melhor dados/variáveis, fontes e período.
\textcolor{gray}{As hipóteses serão examinadas com dados em painel para Mato Grosso. H1 relacionará a conversão da vegetação aos retornos econômicos dos usos alternativos da terra, ao ambiente institucional e à interação entre essas dimensões. H2 relacionará a produtividade agrícola aos choques climáticos, ao estoque prévio de vegetação e à interação entre essas variáveis. As especificações considerarão a ordenação temporal compatível com cada mecanismo, efeitos fixos territoriais e temporais e controles pertinentes a cada hipótese. Dependência espacial, heterogeneidade entre biomas e especificações alternativas serão examinadas por meio de testes e análises de robustez. A interpretação causal dos resultados dependerá das condições de identificação sustentadas pelos dados.}

% ------------------------------------
%%% (VI) DIVISÃO DO TRABALHO
% ------------------------------------

% Falta melhorar: nova divisão do trabalho
\textcolor{gray}{Além desta introdução, o artigo apresenta o contexto econômico do problema e a literatura relacionada aos mecanismos investigados, desenvolve o referencial teórico, descreve a estratégia empírica e os dados, reporta os resultados, discute suas implicações e limitações e, por fim, apresenta as considerações finais.}
